\subsection{Cross-entropy: escape from uniform attention}
\label{app:ce-proof}
Write $u=x_N\bar w\in[-1,1]$, so that $x_N v=s(q,u):=\alpha(q)+(1-\alpha(q))u$ and $\E[u]=-\rho$.
For the original coordinate metric, the exact balanced cross-entropy flow is
\begin{align}
 \ddt{r}&=\frac{r}{2} f(q,r^2),\qquad \ddt{q}=r^2j(q,r^2),\label{eq:ce-flow}\\
 f(q,m)&=\gamma \phi(q)-\E\left[s\tanh\!\left(\frac{m s}{2}\right)\right],\notag\\
 j(q,m)&=\alpha(q)(1-\alpha(q))\,\E\left[(1-u)\left(\gamma-\tanh\!\left(\frac{m s}{2}\right)\right)\right],\notag
\end{align}
with $s=s(q,u)$ inside the expectations.
Here $m$ is a scalar argument; the trained parameters remain $\bm\theta$.
The factor two in the score-gap velocity comes from training both the source score and the shared decoy score.
It is absent if only one gap parameter is trained.
For a two-valued read, $u\in\{-1,1\}$, the factor $1-u$ removes the state $s=1$ and the remaining state has $s=2\alpha-1\le0$ before routing,
so then $j\ge\gamma\phi'(q)$. The anti-copy read takes many values in $[-1,1]$; we only use $|j-\gamma\phi'(q)|\le\tfrac{m}{2}\phi'(q)$,
which keeps $j$ positive while $m<2\gamma$.

\begin{proof}[Proof of \cref{thm:separation}(iv)]
Fix $N,\gamma$ and $\kappa\in\{0,1\}$. Constants below depend on $N,\gamma$, but not on $\varepsilon$.
\proofstep{(a) Exact loss derivatives.}
Write $m=r_1r_2$ inside this proof; then $y z=\xi m s$.
Averaging over the independent error sign, $\E_\xi[\xi\sigma(-\xi a)]=\tfrac{1}{2}(\gamma-\tanh\tfrac{a}{2})$ for the logistic function $\sigma$, so
\begin{align*}
 \pd{\mathcal L_{\mathrm{ce}}}{m}
 &=-\frac{1}{2}\E\left[s\left(\gamma-\tanh\!\left(\frac{m s}{2}\right)\right)\right]
 =-\frac{f(q,m)}{2},\\
 \pd{\mathcal L_{\mathrm{ce}}}{q}
 &=-\frac{m}{2}\E\left[\pd{s}{q}\left(\gamma-\tanh\!\left(\frac{m s}{2}\right)\right)\right]
 =-\frac{m j(q,m)}{2},
\end{align*}
since $\E[s]=\phi(q)$ and $\pd{s}{q}=\alpha(q)(1-\alpha(q))(1-u)$.
The coordinate flow therefore satisfies
\begin{align*}
 \ddt{r_1}&=\frac{r_2f}{2},&\ddt{r_2}&=\frac{r_1f}{2},\\
 \frac{d}{dt}(r_1-r_2)&=-\frac{f}{2}(r_1-r_2),&(r_1-r_2)(0)&=0,\\
 \ddt{q_s}&=\frac{m j}{2},&\ddt{q_d}&=-\frac{m j}{2}.
\end{align*}
Hence $r_1=r_2=r$, $m=r^2$, $q_s=-q_d=\tfrac{q}{2}$, and
\[
 \ddt{r}=\frac{r f}{2},\qquad \ddt{q}=m j,\qquad
 \ddt{m}=m f,\qquad (m(0),q(0))=(\varepsilon^2,0).
\]
\proofstep{(b) Uniform nonlinear bound.}
Using $|\tanh z|\le|z|$, $|s|\le1$, $1-u\ge0$ and $\E[1-u]=1+\rho$,
\begin{align*}
 |f-\gamma \phi|&\le\frac{m}{2}\E[s^2]\le\frac{m}{2},\\
 \left|j-\gamma\phi'(q)\right|&\le\frac{m}{2}\alpha(q)(1-\alpha(q))\E[1-u]=\frac{m}{2}\phi'(q).
\end{align*}
Choose a small fixed $b_0\in(0,\min(1,q_1))$.
Since $\phi'(q)>0$ on its compact closure, there is one constant $C\ge1$ such that
\begin{equation}
 \frac{1}{C}\le j(q,m)\le C,\qquad
 \left|\frac{f(q,m)}{j(q,m)}-\frac{\phi(q)}{\phi'(q)}\right|\le Cm
 \quad(0\le q\le b_0,\ 0\le m\le\gamma).
 \label{eq:ce-ratio-bound}
\end{equation}
Indeed $\tfrac{\gamma}{2}\phi'(q)\le j\le(\gamma+\tfrac{1}{2})\phi'(q)$ for $m\le\gamma$, and the second bound follows from
\[
 \left|\frac{f}{j}-\frac{\phi}{\phi'(q)}\right|
 \le\frac{|f-\gamma \phi|}{j}
 +\left|\frac{\phi}{\phi'(q)}\right|\frac{|\gamma\phi'(q)-j|}{j}.
\]
\proofstep{(c) Score-parametrized comparison.}
For $\kappa=1$, $\tfrac{\phi}{\phi'(q)}=\tfrac{1}{1-\alpha}\in[1,2]$.
For $\kappa=0$, $\phi(0)=0$ and
\[
 \left.\dv{}{q}\left(\frac{\phi(q)}{\phi'(q)}\right)\right|_{q=0}=1.
\]
After decreasing $b_0$, both laws satisfy
\[
 \frac{1}{2}q^{1-\kappa}\le\frac{\phi(q)}{\phi'(q)}\le2q^{1-\kappa}.
\]
For $\kappa=1$, $q^{1-\kappa}$ here means the constant function $1$, including at $q=0$.
While $m\le\gamma$, $j>0$ and \cref{eq:ce-ratio-bound} gives
\[
 \frac{1}{2}q^{1-\kappa}-Cm\le\dv{m}{q}=\frac{f}{j}\le2q^{1-\kappa}+Cm.
\]
Multiplying by the integrating factors $e^{Cq}$ and $e^{-Cq}$ and integrating from zero yields
\begin{align*}
 m(q)&\ge e^{-Cq}\left(\varepsilon^2+\frac{1}{2}\int_0^q e^{Cv}v^{1-\kappa}\,dv\right),\\
 m(q)&\le e^{Cq}\left(\varepsilon^2+2\int_0^q e^{-Cv}v^{1-\kappa}\,dv\right).
\end{align*}
Thus for $0\le q\le b\le b_0$, as long as $m\le\gamma$ on $[0,q]$,
\begin{equation}
 \frac{e^{-Cb}}{4}(\varepsilon^2+q^{2-\kappa})\le m(q)
 \le2e^{Cb}(\varepsilon^2+q^{2-\kappa}).
 \label{eq:ce-polynomial-comparison}
\end{equation}
\proofstep{(d) Existence, gate validity, and time.}
Choose $b\in(0,b_0)$ and $\varepsilon_0\le1$ sufficiently small that
$2e^{Cb}(\varepsilon_0^2+b^{2-\kappa})<\min\{\gamma,4B^2\}$.
By the upper bound in part (c), the first time $m$ could reach $\gamma$ or the gate value $4B^2$ would already have $m<\min\{\gamma,4B^2\}$,
a contradiction, so neither happens before $q=b$; the lower bound prevents exit through $m=0$.
The smooth equation $\dv{m}{q}=\tfrac{f}{j}$ therefore extends to $q=b$.
Define
\[
 T(q)=\int_0^q\frac{dv}{m(v)j(v,m(v))}.
\]
This finite, strictly increasing map reconstructs the initialized network flow:
$\ddt{q}=m j$, $\ddt{m}=m f$, $r=\sqrt m$, $q_s=-q_d=\tfrac{q}{2}$.
Uniqueness and the gate bound identify it with \cref{eq:parameter-flow}.
Combining \cref{eq:ce-ratio-bound,eq:ce-polynomial-comparison} gives
\begin{equation}
 \frac{e^{-Cb}}{2C}\int_0^b\frac{dq}{\varepsilon^2+q^{2-\kappa}}
 \le T(b)\le4Ce^{Cb}\int_0^b\frac{dq}{\varepsilon^2+q^{2-\kappa}}.
 \label{eq:ce-hitting-bounds}
\end{equation}
Finally,
\begin{align*}
 \kappa=1:\quad\int_0^b\frac{dq}{\varepsilon^2+q}
 &=\log\!\left(1+\frac{b}{\varepsilon^2}\right)=\Theta\!\left(\log\!\frac{1}{\varepsilon}\right),\\
 \kappa=0:\quad\int_0^b\frac{dq}{\varepsilon^2+q^2}
 &=\frac{1}{\varepsilon}\arctan\!\left(\frac{b}{\varepsilon}\right)=\Theta\!\left(\frac{1}{\varepsilon}\right).
\end{align*}
Set $q_{\mathrm{esc}}=b$; these are the escape-time orders.\qedhere
\end{proof}

\paragraph{Beyond the escape gap.}
The two-valued argument continued to routing using $j\ge\gamma\phi'(q)$ on all of $[0,q_1]$.
For the anti-copy read this lower bound is only available while $m<2\gamma$, and along the numerical trajectories $m$ exceeds $2\gamma$ before routing.
Numerical integration of \cref{eq:ce-flow} for $N\in\{16,64\}$, $\eta=0.1$, $\kappa\in\{1,0\}$ and $\varepsilon\in[10^{-4},0.03]$ shows that
routing completes in every case, and $j$ stays above $0.04$ ($N=16$) and $0.005$ ($N=64$) along every trajectory.
At $q>q_1$ any positive rule product gives $\operatorname{sign}(z)=b x_N$, as in \cref{app:flow-proof}.

\paragraph{A finite-initialization correction.}
Under $\Dlaw0$, the correlation contribution to the rule gradient is zero at $q=0$.
For $r>0$, cross-entropy instead gives $f(0,r^2)<0$, with expansion
\[
 f(0,r^2)=-\frac{r^2}{2}\E[v(0)^2]+O(r^6).
\]
Thus $\ddt{r}$ initially shrinks at order $r^3$, even though the leading credit cancels.
This correction is compatible with the routing theorem's local comparison: it changes the comparison by $O(r^2)$ without removing the quadratic growth in the score gap.
